\documentclass[10pt,a4paper]{amsart}
\usepackage{amsmath,amssymb,amsthm,graphicx}
\usepackage{verbatim}
\newenvironment{reference}{\comment}{\endcomment}
%\newenvironment{reference}{}{}
\newenvironment{slogan}{\comment}{\endcomment}
\newenvironment{history}{\comment}{\endcomment}
% For commutative diagrams we use Xy-pic
\usepackage[all]{xy}

% We use 2cell for 2-commutative diagrams.
\xyoption{2cell}
% Theorem environments.
%
\theoremstyle{plain}
\newtheorem{theorem}[subsection]{Theorem}
\newtheorem{proposition}[subsection]{Proposition}
\newtheorem{lemma}[subsection]{Lemma}

\theoremstyle{definition}
\newtheorem{definition}[subsection]{Definition}
\newtheorem{example}[subsection]{Example}
\newtheorem{exercise}[subsection]{Exercise}
\newtheorem{situation}[subsection]{Situation}

\theoremstyle{remark}
\newtheorem{remark}[subsection]{Remark}
\newtheorem{remarks}[subsection]{Remarks}

\numberwithin{equation}{subsection}

% Macros
%
\def\lim{\mathop{\mathrm{lim}}\nolimits}
\def\colim{\mathop{\mathrm{colim}}\nolimits}
\def\Spec{\mathop{\mathrm{Spec}}}
\def\Hom{\mathop{\mathrm{Hom}}\nolimits}
\def\Ext{\mathop{\mathrm{Ext}}\nolimits}
\def\SheafHom{\mathop{\mathcal{H}\!\mathit{om}}\nolimits}
\def\SheafExt{\mathop{\mathcal{E}\!\mathit{xt}}\nolimits}
\def\Sch{\mathit{Sch}}
\def\Mor{\mathop{\mathrm{Mor}}\nolimits}
\def\Ob{\mathop{\mathrm{Ob}}\nolimits}
\def\Sh{\mathop{\mathit{Sh}}\nolimits}
\def\NL{\mathop{N\!L}\nolimits}
\def\CH{\mathop{\mathrm{CH}}\nolimits}
\def\proetale{{pro\text{-}\acute{e}tale}}
\def\etale{{\acute{e}tale}}
\def\QCoh{\mathit{QCoh}}
\def\Ker{\mathop{\mathrm{Ker}}}
\def\Im{\mathop{\mathrm{Im}}}
\def\Coker{\mathop{\mathrm{Coker}}}
\def\Coim{\mathop{\mathrm{Coim}}}

% Boxtimes
%
\DeclareMathSymbol{\boxtimes}{\mathbin}{AMSa}{"02}

%
% Macros for moduli stacks/spaces
%
\def\QCohstack{\mathcal{QC}\!\mathit{oh}}
\def\Cohstack{\mathcal{C}\!\mathit{oh}}
\def\Spacesstack{\mathcal{S}\!\mathit{paces}}
\def\Quotfunctor{\mathrm{Quot}}
\def\Hilbfunctor{\mathrm{Hilb}}
\def\Curvesstack{\mathcal{C}\!\mathit{urves}}
\def\Polarizedstack{\mathcal{P}\!\mathit{olarized}}
\def\Complexesstack{\mathcal{C}\!\mathit{omplexes}}
% \Pic is the operator that assigns to X its picard group, usage \Pic(X)
% \Picardstack_{X/B} denotes the Picard stack of X over B
% \Picardfunctor_{X/B} denotes the Picard functor of X over B
\def\Pic{\mathop{\mathrm{Pic}}\nolimits}
\def\Picardstack{\mathcal{P}\!\mathit{ic}}
\def\Picardfunctor{\mathrm{Pic}}
\def\Deformationcategory{\mathcal{D}\!\mathit{ef}}
\usepackage{cleveref}
\crefname{originallemma}{original lemma}{original lemmas}
\Crefname{originallemma}{Original lemma}{Original lemmas}
\makeatletter
\newlabel{lemma-sitting-in-degrees}{{lemma-sitting-in-degrees}{source}{}{sourcepointer-1}{}}
\newlabel{lemma-sitting-in-degrees@cref}{{[originallemma][0][]lemma-sitting-in-degrees}{[1][0][]source}{}{}{}}
\makeatother
\crefname{originallemma}{original lemma}{original lemmas}
\Crefname{originallemma}{Original lemma}{Original lemmas}
\makeatletter
\newlabel{lemma-top-coh-divisible}{{lemma-top-coh-divisible}{source}{}{sourcepointer-2}{}}
\newlabel{lemma-top-coh-divisible@cref}{{[originallemma][0][]lemma-top-coh-divisible}{[1][0][]source}{}{}{}}
\makeatother
\crefname{originallemma}{original lemma}{original lemmas}
\Crefname{originallemma}{Original lemma}{Original lemmas}
\makeatletter
\newlabel{lemma-cd-top-vanishing-some-module}{{lemma-cd-top-vanishing-some-module}{source}{}{sourcepointer-3}{}}
\newlabel{lemma-cd-top-vanishing-some-module@cref}{{[originallemma][0][]lemma-cd-top-vanishing-some-module}{[1][0][]source}{}{}{}}
\makeatother
\crefname{originalreference}{original reference}{original references}
\Crefname{originalreference}{Original reference}{Original references}
\makeatletter
\newlabel{algebra-lemma-silly}{{algebra-lemma-silly}{source}{}{sourcepointer-4}{}}
\newlabel{algebra-lemma-silly@cref}{{[originalreference][0][]algebra-lemma-silly}{[1][0][]source}{}{}{}}
\makeatother
\crefname{originalreference}{original reference}{original references}
\Crefname{originalreference}{Original reference}{Original references}
\makeatletter
\newlabel{derived-remark-truncation-distinguished-triangle}{{derived-remark-truncation-distinguished-triangle}{source}{}{sourcepointer-5}{}}
\newlabel{derived-remark-truncation-distinguished-triangle@cref}{{[originalreference][0][]derived-remark-truncation-distinguished-triangle}{[1][0][]source}{}{}{}}
\makeatother
\crefname{originalreference}{original reference}{original references}
\Crefname{originalreference}{Original reference}{Original references}
\makeatletter
\newlabel{derived-lemma-trick-vanishing-composition}{{derived-lemma-trick-vanishing-composition}{source}{}{sourcepointer-6}{}}
\newlabel{derived-lemma-trick-vanishing-composition@cref}{{[originalreference][0][]derived-lemma-trick-vanishing-composition}{[1][0][]source}{}{}{}}
\makeatother
\crefname{originalreference}{original reference}{original references}
\Crefname{originalreference}{Original reference}{Original references}
\makeatletter
\newlabel{dualizing-lemma-nonvanishing-generically-local}{{dualizing-lemma-nonvanishing-generically-local}{source}{}{sourcepointer-7}{}}
\newlabel{dualizing-lemma-nonvanishing-generically-local@cref}{{[originalreference][0][]dualizing-lemma-nonvanishing-generically-local}{[1][0][]source}{}{}{}}
\makeatother
\title{Vanishing of the zeroth local cohomology of the dualizing complex forces the cohomological dimension below the dimension}
\author{The Stacks Project Authors}
\date{stacks-project:local-cohomology.tex:local-snapshot-20261012; GNU FDL 1.2 or later}
\begin{document}
\maketitle
\noindent\textit{Editorial note.} All delivered mathematical text is the author's own, copied from the pinned Stacks Project chapter source sources/stacks-local-cohomology/pinned-original-source.tex; no statement or proof step was rewritten, reordered or removed.\par\medskip
\noindent\textit{Editorial note.} Preamble provenance: the Stacks Project keeps one single shared preamble for every chapter; this card takes it directly from the staged sources/stacks-local-cohomology/preamble.tex as three verbatim spans (lines 8-12, 14-18 and 154-217), so no derived file was written and nothing was added or redefined. The wrapper class and the base packages are supplied by the assembler.\par\medskip
\noindent\textit{Editorial note.} Bibliography provenance: neither the counted statement nor its proof contains a citation (this lane rescans the whole span before writing the card), so the card carries no bibliography and no entry was reconstructed.\par\medskip
\noindent\textit{Editorial note.} Selection provenance: this card was opened under the file-based control contract (CONTROL.md), with the mechanical pre-card gate run on its statement and proof interval before the card was written.\par\medskip
\noindent\textit{Allowed standard prerequisites (editorial).}\begin{itemize}
\item The degrees in which the cohomology of a normalized dualizing complex is nonzero and the bound on the dimension of the support of each of them. Invoked early in the counted proof (source lines 3659-3663) as a pointer to a lemma of this chapter. Identifies the top cohomology module of the complex and lets the proof kill the lower ones by a single element.
\item Prime avoidance, the truncation triangle of a complex, and the vanishing of an endomorphism of a truncation below the truncation degree. Invoked inside the counted proof (source lines 3665-3677) as cross-chapter pointers. Produces an element outside the minimal primes of maximal dimension that kills the truncated part, so that a power of it factors through the top cohomology module.
\end{itemize}\medskip
\ifdefined\hypertarget\hypertarget{sourcepointer-1}{}\fi
\noindent\textit{Source reference.} Original lemma lemma-sitting-in-degrees: sources/stacks-local-cohomology/pinned-original-source.tex, source lines 3663--3663. Navigation pointer only: the counted text refers to the Stacks item lemma-sitting-in-degrees, which belongs to another chapter of the project and is not part of this card. The pointer carries the target's own label as its number because this chapter's staged snapshot does not contain the other chapter's tag table; its content is not reproduced here.\par\medskip
\ifdefined\hypertarget\hypertarget{sourcepointer-2}{}\fi
\noindent\textit{Source reference.} Original lemma lemma-top-coh-divisible: sources/stacks-local-cohomology/pinned-original-source.tex, source lines 3685--3685. Navigation pointer only: the counted text refers to the Stacks item lemma-top-coh-divisible, which belongs to another chapter of the project and is not part of this card. The pointer carries the target's own label as its number because this chapter's staged snapshot does not contain the other chapter's tag table; its content is not reproduced here.\par\medskip
\ifdefined\hypertarget\hypertarget{sourcepointer-3}{}\fi
\noindent\textit{Source reference.} Original lemma lemma-cd-top-vanishing-some-module: sources/stacks-local-cohomology/pinned-original-source.tex, source lines 3686--3686. Navigation pointer only: the counted text refers to the Stacks item lemma-cd-top-vanishing-some-module, which belongs to another chapter of the project and is not part of this card. The pointer carries the target's own label as its number because this chapter's staged snapshot does not contain the other chapter's tag table; its content is not reproduced here.\par\medskip
\ifdefined\hypertarget\hypertarget{sourcepointer-4}{}\fi
\noindent\textit{Source reference.} Original reference algebra-lemma-silly: sources/stacks-local-cohomology/pinned-original-source.tex, source lines 1466--1466. Navigation pointer only: the counted text refers to the Stacks item algebra-lemma-silly, which belongs to another chapter of the project and is not part of this card. The pointer carries the target's own label as its number because this chapter's staged snapshot does not contain the other chapter's tag table; its content is not reproduced here.\par\medskip
\ifdefined\hypertarget\hypertarget{sourcepointer-5}{}\fi
\noindent\textit{Source reference.} Original reference derived-remark-truncation-distinguished-triangle: sources/stacks-local-cohomology/pinned-original-source.tex, source lines 3674--3674. Navigation pointer only: the counted text refers to the Stacks item derived-remark-truncation-distinguished-triangle, which belongs to another chapter of the project and is not part of this card. The pointer carries the target's own label as its number because this chapter's staged snapshot does not contain the other chapter's tag table; its content is not reproduced here.\par\medskip
\ifdefined\hypertarget\hypertarget{sourcepointer-6}{}\fi
\noindent\textit{Source reference.} Original reference derived-lemma-trick-vanishing-composition: sources/stacks-local-cohomology/pinned-original-source.tex, source lines 3675--3675. Navigation pointer only: the counted text refers to the Stacks item derived-lemma-trick-vanishing-composition, which belongs to another chapter of the project and is not part of this card. The pointer carries the target's own label as its number because this chapter's staged snapshot does not contain the other chapter's tag table; its content is not reproduced here.\par\medskip
\ifdefined\hypertarget\hypertarget{sourcepointer-7}{}\fi
\noindent\textit{Source reference.} Original reference dualizing-lemma-nonvanishing-generically-local: sources/stacks-local-cohomology/pinned-original-source.tex, source lines 849--849. Navigation pointer only: the counted text refers to the Stacks item dualizing-lemma-nonvanishing-generically-local, which belongs to another chapter of the project and is not part of this card. The pointer carries the target's own label as its number because this chapter's staged snapshot does not contain the other chapter's tag table; its content is not reproduced here.\par\medskip
\begin{lemma}
\label{lemma-cd-bound-dualizing}
Let $A$ be a Noetherian local ring with
normalized dualizing complex $\omega_A^\bullet$.
Let $I \subset A$ be an ideal.
If $H^0_{V(I)}(\omega_A^\bullet) = 0$, then $\text{cd}(A, I) < \dim(A)$.
\end{lemma}

\begin{proof}
Set $d = \dim(A)$. Let $\mathfrak p_1, \ldots, \mathfrak p_n \subset A$
be the minimal primes of dimension $d$.
Recall that the finite $A$-module
$H^{-i}(\omega_A^\bullet)$ is nonzero only for
$i \in \{0, \ldots, d\}$ and that the support
of $H^{-i}(\omega_A^\bullet)$ has dimension $\leq i$, see
Lemma \ref{lemma-sitting-in-degrees}.
Set $\omega_A = H^{-d}(\omega_A^\bullet)$.
By prime avoidance (Algebra, Lemma \ref{algebra-lemma-silly})
we can find $f \in A$, $f \not \in \mathfrak p_i$
which annihilates $H^{-i}(\omega_A^\bullet)$ for $i < d$.
Consider the distinguished triangle
$$
\omega_A[d] \to \omega_A^\bullet \to
\tau_{\geq -d + 1}\omega_A^\bullet \to \omega_A[d + 1]
$$
See Derived Categories, Remark
\ref{derived-remark-truncation-distinguished-triangle}.
By Derived Categories, Lemma \ref{derived-lemma-trick-vanishing-composition}
we see that $f^d$ induces the zero endomorphism of
$\tau_{\geq -d + 1}\omega_A^\bullet$.
Using the axioms of a triangulated category, we find a map
$$
\omega_A^\bullet \to \omega_A[d]
$$
whose composition with $\omega_A[d] \to \omega_A^\bullet$ is
multiplication by $f^d$ on $\omega_A[d]$.
Thus we conclude that $f^d$ annihilates $H^d_{V(I)}(\omega_A)$.
By Lemma \ref{lemma-top-coh-divisible} we conlude $H^d_{V(I)}(\omega_A) = 0$.
Then we conclude by Lemma \ref{lemma-cd-top-vanishing-some-module}
and the fact that $(\omega_A)_{\mathfrak p_i}$ is nonzero
(see for example
Dualizing Complexes, Lemma
\ref{dualizing-lemma-nonvanishing-generically-local}).
\end{proof}

\end{document}
